% Direct-authority diplomatic French for printed p.153.
\SourcePage{153}

\LayerTarget{tohoku:construction:four-functor-square:p2}
\noindent objets $G$-acycliques.

\LayerTarget{tohoku:prop:2.5.5}
\noindent\textsc{Proposition 2.5.5.} --- \textit{Plaçons-nous dans les
conditions précédentes. Pour un $A\in C$, considérons les deux suites
spectrales de hyperhomologie de $G$ par rapport au complexe $\mathbf F(A)$,
ce sont des foncteur spectraux en $A$. Ils sont isomorphes respectivement au
deuxième foncteur spectral du foncteur composé $F'G'$ et $GF$.}

\LayerTarget{tohoku:prop:2.5.5:proof}
L'assertion relative au foncteur composé $GF$ n'est autre que la proposition
2.5.4. Quant à la deuxième suite spectrale du foncteur composé $F'G'$, c'est
par définition la deuxième suite spectrale de $F'$ par rapport au complexe
$G'(C(A))$ (où $C(A)$ désigne une résolution injective de $A$), donc
s'identifie en vertu de proposition 2.5.3 à la deuxième suite spectrale du
bicomplexe $\mathbf F'G'(C(A))$ (dont le premier degré est celui qui provient
de $C(A)$), or on a $\mathbf F'G'=G\mathbf F$, il faut donc calculer la suite
spectrale $\mathrm{II}(G\mathbf F(C(A)))$. Or pour tout $q$ fixé,
$\mathbf F'^qC(A)$ est une résolution de $\mathbf F'^q(A)$ (puisque
$\mathbf F'^q$ est exact) par des objets $G$-acycliques, donc
$\mathbf F'(C(A))$ devient, après échange de ses deux degrés, une résolution
de $\mathbf F'(A)$ par des objets $G$-acycliques. En vertu de 2.4, il
s'ensuit que la deuxième suite spectrale de $G(\mathbf F'(C(A)))$ s'identifie
à la première suite spectrale d'hyperhomologie de $G$ par rapport au complexe
$\mathbf F'(A)$.

\LayerTarget{tohoku:cor:2.5.5}
\noindent\textsc{Corollaire.} --- \textit{$C,C',C'',C'''$ et $F,G,F',G'$
étant comme ci-dessus (voir diagramme) on suppose donnés des foncteurs
résolvants $\mathbf F,\mathbf G,\mathbf F',\mathbf G'$ pour $F,G,F',G'$, on
suppose $\mathbf F'\mathbf G'=\mathbf G\mathbf F$ (isomorphismes fonctories
compatibles avec les opérateurs de dérivation dans les foncteurs résolvants),
et que $\mathbf F$ (resp. $\mathbf G'$) transforme objets injectifs en objets qui sont
$G$-acycliques (resp. $F'$-acycliques). Alors pour tout $A\in C$, les
deuxièmes suites spectrales des foncteurs composés $F'G'$ et $GF$ pour $A$
s'identifient à la première et deuxième suite spectrale d'hyperhomologie de
$G$ pour le complexe $\mathbf F(A)$, ou encore à la première et la deuxième
suite spectrale d'hyperhomologie de $F'$ pour le complexe $\mathbf G'(A)$, ou
enfin à l'une et l'autre des deux suites spectrales du bicomplexe
$\mathbf G(\mathbf F(A))=\mathbf F'(\mathbf G'(A))$.}

\LayerTarget{tohoku:chapter:III}
\noindent\textbf{Chapitre III. Cohomologie a coefficients dans un faisceau.}

\LayerTarget{tohoku:section:3.1}
\noindent\textbf{3.1. Généralités sur les faisceaux.}

\LayerTarget{tohoku:def:presheaf-and-sheaf:p1}
Soit $X$ un espace topologique (non nécessairement séparé). Rappelons (1.7
exemple h) qu'on appelle \textit{préfaisceau d'ensembles sur $X$} tout système
inductif d'ensembles défini sur l'ensemble des parties ouvertes non vides de
$X$, ordonné par $\supset$. Il consiste donc en la donnée, pour tout ouvert
$U\subset X$, d'un ensemble $F(U)$, et pour tout couple d'ouverts non vides
$U,V$ avec $U\supset V$, d'une \textit{application de restrication}
\[
\varphi_{VU}:F(U)\longrightarrow F(V),
\]
avec les conditions : $\varphi_{UU}=$ application identique de $F(U)$,
$\varphi_{WV}\varphi_{VU}=\varphi_{WU}$ si
$U\supset V\supset W$. On dit que le préfaisceau $F$ est un \textit{faisceau}
si pour tout recouvrement $(U_i)$ d'un ouvert $U$ de $X$ par des ouverts non
vides, et toute famille $(f_i)$ d'éléments $f_i\in F(U_i)$ telle que
\[
\varphi_{U_{ij}U_i}f_i=\varphi_{U_{ij}U_j}f_j
\]
pour tout couple $(i,j)$ tel que $U_{ij}=U_i\cap U_j\ne\varnothing$, il existe
un $f\in F(U)$ et un seul tel que $\varphi_{U_iU}f=f_i$ pour tout $i$. Si dans
les définitions précédentes, on suppose

\TohokuContinuation{154}
